\section{Discussion}
\label{sec:discussion}

The acquisition-contract view connects an algorithmic privacy question to Web
architecture. A failure at T1 implicates image-serving infrastructure; a failure
at T3 implicates what an agent observes after browser rendering; a failure at
T5 or T6 shows that a pixel-only refusal cannot survive generative mediation.
Independent layers make those diagnoses possible, while a cascade would not.

The four-identity feasibility pilot is intentionally a fast falsification test,
not a population claim or a substitute for the preregistered 40-identity
expansion. It is far too small for confidence-heavy ranking or demographic
subgroup conclusions.
Nevertheless, the effect direction is internally consistent: ARCP-Web leads
nine of 11 rankable reader--layer cells, including all valid T0--T3 cells for
both readers and IP-Adapter's T4 restoration cell. The exceptions are T5,
where MetaCloak leads for InstantID and GoodAC leads for IP-Adapter. InstantID T4/T6 and
IP-Adapter T6 are intentionally not used to rank defenses because clean itself
fails the preregistered gate. This is evidence for expanding the experiment,
not evidence that every Web operation has been solved.
Real platform upload/download traces, additional identity distributions, and
adaptive attacks remain necessary. ARCP-Web's training proxies approximate T4--T6
rather than differentiating through their exact evaluation implementations, so
transfer is an empirical risk. Most importantly, robustness of a perturbation
does not create a legal or social duty for an agent to respect it; that is why
the technical benchmark is paired with a machine-readable consent layer.
